% Generated draft: requires mathematical and source review.
\documentclass[11pt]{article}
\usepackage[margin=25mm]{geometry}
\usepackage{fontspec}
\setmainfont{Latin Modern Roman}
\usepackage{amsmath,amssymb,mathtools,unicode-math}
\setmathfont{Latin Modern Math}
\usepackage{enumitem,needspace,xurl,hyperref}
\setlength{\parindent}{0pt}
\setlength{\parskip}{5pt}
\setlength{\emergencystretch}{3em}
\newcommand{\N}{\mathbb{N}}
\newcommand{\Z}{\mathbb{Z}}
\newcommand{\Q}{\mathbb{Q}}
\newcommand{\R}{\mathbb{R}}
\newcommand{\C}{\mathbb{C}}
\newcommand{\F}{\mathbb{F}}
\newcommand{\A}{\mathbb{A}}
\newcommand{\PP}{\mathbb{P}}
\newcommand{\E}{\mathbb{E}}
\newcommand{\eps}{\varepsilon}
\newcommand{\dd}{\,\mathrm{d}}
\begin{document}
\section*{UnsolvedMath ID: 700598}
\textbf{Source problem code:} catalogue entry
\subsection{Definitions and mathematical statement}
A group \(H\) is \emph{periodic} if every element of \(H\) has finite order. Define its derived series by \(H^{(0)}=H\) and \(H^{(n+1)}=[H^{(n)},H^{(n)}]\), where \([K,K]\) is the subgroup generated by all commutators of elements of \(K\). The group \(H\) is \emph{soluble} if \(H^{(n)}=\{1\}\) for some integer \(n\geq 0\). A group is \emph{locally soluble} if every finitely generated subgroup is soluble; no uniform bound on their derived lengths is assumed.

For subgroups \(A,B\leq G\), write \(AB=\{ab:a\in A,\ b\in B\}\). Thus \(G=AB\) means that every element of \(G\) is a product of an element of \(A\) and an element of \(B\); neither subgroup is assumed normal.

\textbf{Question (Ya. P. Sysak).} For every locally soluble group \(G\) and all periodic subgroups \(A,B\leq G\) satisfying \(G=AB\), must \(G\) be periodic? Explicitly, do these hypotheses imply
\[
\forall g\in G\;\exists n\in\Z_{\geq 1}\quad g^n=1?
\]
\subsection{Short English statement}
Must every locally soluble group that is a product of two periodic subgroups itself be periodic? Here periodic means that every element has finite order, and the product is the setwise product of the two subgroups.
\subsection{Sources}
\begin{itemize}
\item Supplied local UnsolvedMath registry record: permanent ID \(700598\), exact source code \texttt{catalogue entry}, title \emph{Kourovka Notebook Problem 9.65}.
\item The supplied record attributes the problem to Ya. P. Sysak (1984), \emph{The Kourovka Notebook (Unsolved Problems in Group Theory)}, 9th issue (1984), original problem number 9.65. It also identifies the 21st issue edition (2026), arXiv:1401.0300v48. The supplied collection description names E. I. Khukhro and V. D. Mazurov as editors.
\item The registry labels the problem \emph{open} and records an update timestamp of \texttt{2026-10-16T14:00:00Z}. This is supplied registry metadata, not an independently verified mathematical or literature assessment.
\end{itemize}
\end{document}
